% Created 2022-06-06 Mon 07:47
% Intended LaTeX compiler: pdflatex
\documentclass[11pt,a4paper]{article}
\usepackage[utf8]{inputenc}
\usepackage[T1]{fontenc}
\usepackage{graphicx}
\usepackage{longtable}
\usepackage{wrapfig}
\usepackage{rotating}
\usepackage[normalem]{ulem}
\usepackage{amsmath}
\usepackage{amssymb}
\usepackage{capt-of}
\usepackage{hyperref}
\usepackage[type={CC},modifier={by-nc-sa},version={3.0}]{doclicense}
\usepackage[top=1cm,bottom=1cm,left=1cm,right=1cm]{geometry}
\usepackage{pdflscape}
\usepackage[final]{pdfpages}
\usepackage[myheadings]{fullpage}
\usepackage{enumerate}
\usepackage{parskip}
\usepackage{dashrule}
\usepackage{enumitem}
\usepackage{sectsty}
\usepackage{latexsym}
\usepackage{framed}
\usepackage{pbox}
\usepackage{ifdraft}
\usepackage[english]{babel}
\usepackage{fourier}
\usepackage{shadowtext}
\usepackage[utf8]{inputenc}
\usepackage[T1]{fontenc}
\usepackage[protrusion=true, expansion=true]{microtype}
\usepackage{textcomp}
\usepackage{soul}
\usepackage{lipsum}
\usepackage{blindtext}
\usepackage{graphicx}
\usepackage{float}
\usepackage{wrapfig}
\usepackage{rotating}
\usepackage{tikz}
\usepackage{tikzpagenodes}
\usepackage[font=small, labelfont=bf]{caption}
\usepackage{pgfgantt}
\usepackage{amsmath}
\usepackage{bm}
\usepackage{amssymb}
\usepackage{marvosym}
\usepackage{mathpazo}
\usepackage{mathtools}
\usepackage{siunitx}
\usepackage{tabularx}
\usepackage{longtable}
\usepackage{array}
\usepackage{makecell}
\usepackage{totpages,fancyhdr}
\def\ORGTITLE  {Literature Review}
\def\ORGAUTHOR {DP@H_2O}
\pagestyle{fancy}
\fancyhf{}
\setlength\headheight{1cm}
\fancyhead[L]{\LARGE \ORGTITLE  }
\fancyhead[R]{\large \ORGAUTHOR }
\fancyfoot[C]{\small Page \thepage}
\usepackage[natbib,style=authoryear-comp]{biblatex}
\bibliography{../references/refs_bbt.bib}
\usepackage{url}
\usepackage{booktabs}
\usepackage{hyperref}
\hypersetup{colorlinks=true,linktoc=all,linkcolor=blue,filecolor=magenta,urlcolor=cyan,citecolor=magenta}
\usepackage{setspace}
\usepackage{titling}
\title{\textsc{Research Plan}}
\author{\href{mailto:daniel.atwater@utas.edu.au}{Daniel Patrick Atwater}}
\setcounter{secnumdepth}{1}
\setcounter{tocdepth}{1}
\usepackage{xcolor}
\newcommand{\crule}[3][black]{\textcolor{#1}{\rule{#2}{#3}}}
\newcommand{\source}[1]{\caption*{\hfill Source: {#1}} }
\newcommand{\HRule}[1]{\rule{\linewidth}{#1}}
\newcommand{\latlonDec}[2]{$\ang{#1}$ \textbf{#2}}
\newcommand{\citationNeeded}{\textbf{}}
\newcommand{\htwoo}{H$_2$O }
\newcommand{\otwo}{O$_2$ }
\newcommand{\cotwo}{CO$_2$ }
\newcommand{\chfour}{CH$_4$ }
\newcommand{\ntwo}{N$_2$ }
\newcommand{\moc}{\textit{\textsc{moc }}} %Meridional Overturning Circulation
\newcommand{\nadw}{\textit{\textsc{nadw }}} %Meridional Overturning Circulation
\newcommand{\sip}{\textit{\textsc{sip }}} %Sea ice production
\newcommand{\dsw}{\textit{\textsc{{dsw }}}} %Dense Shelf Water
\newcommand{\aabw}{\textit{\textsc{{aabw }}}} %Antarctic Bottom Water
\newcommand{\asf}{\textit{\textsc{{asf }}}} %Antarctic Slope Front
\newcommand{\asc}{\textit{\textsc{{asc }}}} %Antarctic Slope Current
\newcommand{\acc}{\textit{\textsc{{acc }}}} %Antarctic Slope Current
\newcommand{\cdp}{\textit{\textsc{{cdp }}}} %Cape Darnley Polynya
\newcommand{\cdfi}{\textit{\textsc{{cdfi }}}}
\newcommand{\mc}{\textit{\textsc{{mc }}}} %Mawson Coast
\newcommand{\ais}{\textit{\textsc{{ais} }}} %Amery Ice Shelf
\newcommand{\fastice}{fast ice }
\newcommand{\Fastice}{Fast ice }
\newcommand{\Lpolynya}{$\mathcal{L}$-polynya }
\newcommand{\Tpolynya}{$\mathcal{T}$-polynya }
\newcommand{\wmo}{\href{https://public.wmo.int/en}{\uppercase{wmo}}}
\newcommand{\utas}{\href{https://utas.edu.au}{\textsc{UTas}}}
\newcommand{\aapp}{\href{https://aappartnership.org.au}{\uppercase{aapp}}}
\newcommand{\imas}{\href{https://www.imas.utas.edu.au}{\uppercase{imas}}}
\newcommand{\csiro}{\href{https://www.csiro.au/en/about/people/business-units/oceans-and-atmosphere}{\uppercase{csiro}}}
\newcommand{\bom}{\href{http://www.bom.gov.au}{\uppercase{bom}}}
\newcommand{\anu}{\href{http://www.anu.edu.au}{\uppercase{anu}}}
\newcommand{\primarysupervisor}{\href{mailto:alexander.fraser@utas.edu.au}{\textsc{Alex Fraser}}}
\newcommand{\secondarysupervisor}{\href{mailto:will.hobbs@utas.edu.au}{\textsc{Will Hobbs}}}
\newcommand{\tertiarysupervisor}{\href{mailto:paul.spence@utas.edu.au}{\textsc{Pual Spence}}}
\newcommand{\grc}{\href{mailto:maxim.nikurashin@utas.edu.au}{\textsc{Max Nikurashin}}}
\author{DP@H\textsubscript{2}O}
\date{\textit{<2022-06-01 Wed>}}
\title{Australian Fast Ice Model\\\medskip
\large Literature Review}
\hypersetup{
 pdfauthor={DP@H\textsubscript{2}O},
 pdftitle={Australian Fast Ice Model},
 pdfkeywords={},
 pdfsubject={},
 pdfcreator={Emacs 28.1 (Org mode 9.5.2)}, 
 pdflang={English}}
\begin{document}

\maketitle
\tableofcontents


\section{Introduction}
\label{sec:org8a33e7c}

The meridional overturning circulation (\moc) within the Earth's ocean
basins is understood to play a significant contribution in regulating
heat content of the climate, ventilating the deep ocean and \cotwo uptake
\citep{orsiCirculationMixingProduction1999,levitusWarmingWorldOcean2005,bindoffObservationsOceanicClimate2007,purkeyWarmingGlobalAbyssal2010,talleyClosureGlobalOverturning2013}. With
overturning time scales on the order of a millennia \citep{johnsonQuantifyingAntarcticBottom2008}
understanding the driving mechanisms of the \moc has implications for
the rates of change of key parameters in Earth's climate: e.g. \cotwo uptake rate, latent heat draw down, and ventilation of the deep ocean
\citep{bindoffObservationsOceanicClimate2007,purkeyWarmingGlobalAbyssal2010}. 

The conceptual model of the \moc includes the tropical waters of the
Atlantic flowing north via the Gulf Stream and overturning at various
locations in the north Atlantic due to baroclinic instabilities, thus
becoming North Atlantic Deep Water. At the other end of planet, along
the coast of Antarctica, a complex set of seasonal events give rise to
coastal processes that cause cold dense water to form at the surface
at key locations \citep{orsiCirculationMixingProduction1999}. As the water forms it is
highly unstable and begins to sink flowing down the slopes of the
continental shelf, mixing with ambient waters in its descent, and in
the process becoming dense shelf water (\dsw) \citep{orsiCirculationMixingProduction1999}. Both
Ekman forced and baroclinic instabilities transport
\citep{orsiCirculationMixingProduction1999} this \dsw flows across the Antarctic slope front,
with some indication of mixing
\citep{mizutaSeasonalEvolutionCape2021,vanachterModellingLandfastSea2022,rintoulOceanHeatDrives2016}, and into the various
abyssal plains (e.g. becoming Antarctic Bottom Water \textemdash
\aabw) \citep{stewartAccelerationOverturningAntarctic2019}. Abyssal
waters with a southern origin that are colder than the North Atlantic
Deep Water, with temperatures below \si{2}{\degress\Celisus} have been referred to
generically as Antarctic Bottom Water (\aabw) for quite a long time [Deacon,
1933; Speer \& Zenk, 1993]. This water mass fills the ocean basins of
the Southern Hemisphere \citep{johnsonQuantifyingAntarcticBottom2008}, and is identified as cold
dense abyssal water originating from overturning coastal waters around
Antarctica \citep{orsiCirculationMixingProduction1999}. 

In this manner, \aabw mass accounts for 30\textendash 40\% of global ocean
mass \citep{johnsonQuantifyingAntarcticBottom2008}. Its formation and controlling mechanisms are
areas of contemporary research. In the effort to prognosticate furture
climate, the current capability to constrain the climate's heat
budget, via global climate models, is dependent upon knowledge of
these controlling mechanisms of the deep ocean
\citep{dominguesImprovedEstimatesUpperocean2008,murphyObservationallyBasedEnergy2009,purkeyGlobalContractionAntarctic2012,durackOceanWarmingSurface2018}. At
present, the uncertainty in deep-ocean heat uptake has been shown to be a
significant cause of variation among climate models
\citep{boeDeepOceanHeat2009,durackOceanWarmingSurface2018}. While on the other hand, empirically, it
also has been shown that Earth's systems linked to the long-term
water-phase-mass distribution and surface temperature variations are
accelerating \citep{murphyObservationallyBasedEnergy2009,portnerIPCCSpecialReport2019}.

One specific uncertainty of the \moc under investigation in this work,
is the rate at which \aabw is formed via \dsw formation. To
investigate this uncertainty I will be undertaking a series of
simulations, and likely enhancements/tuning of those simulations, of a
coupled physical system (sea ice-ocean and possibly an atmosphere
coupled model) across the circumpolar Antarctic domain.

\subsection{Project Overview}
\label{sec:orgf57ca88}

First, before analysing and revising the physically simulated
environment across the whole of the Antarctic coastline for fast ice
distribution accuracy, I will be focusing on specific region of East
Antarctica. The Mawson Coast and Prydz Bay region of East Antarctica under investigation here, is chosen due to its:
\begin{enumerate}
\item significance in \aabw formation
  \citep{williamsSurfaceOceanographyBROKEWest2010,ohshimaGlobalViewSeaice2016,tamuraSeaIceProduction2016,williamsg.d.SuppressionAntarcticBottom,fraserLandfastIceControls2019,aokiFresheningAntarcticBottom2020},
\item previous remote sensing, oceanographic a geographic studies (i.e. data availability)
  \citep{meijersCirculationWaterMasses2010,williamsSurfaceOceanographyBROKEWest2010,ohshimaAntarcticBottomWater2013,tamuraSeaIceProduction2016,fraserLandfastIceControls2019,aokiFresheningAntarcticBottom2020,smithCapeDarnleyBathymetry2020},
  and
\item previous geophysical fluid modelling covering or adjacent to this region
  \citep{galton-fenziModelingBasalMelting2012,nakayamaNumericalInvestigationFormation2014,liuModelingModifiedCircumpolar2017,kusaharaModelingOceanCryosphere2017,kusaharaDenseShelfWater2017,naughtenIntercomparisonAntarcticIceshelf2018,vanachterModellingLandfastSea2022}.
\end{enumerate}

The distinction of the Mawson Coast, in this context, is its inclusion of the Cape Darnley Polynya
(\cdp) \citep{meijersCirculationWaterMasses2010}. The \cdp is a coastal ocean polynya
(i.e. latent\textendash heat polynya) \citep{massomDistributionFormativeProcesses1998,tamuraMappingSeaIce2008}. Some coastal ocean
polynyas are sources of \aabw \citep{massomDistributionFormativeProcesses1998,maquedam.PolynyaDynamicsReview}. For
reasons of clarification, I will use the \cite{barberChapterRoleSea2007} definition of
polynyas:
\begin{quote}
  Polynyas are persistent reoccurring regions of open water and/or thin ice,
  tens to tens of thousands of square kilometres in extent, that occur at
  locations where more consolidated and thicker ice cover would be
  climatologically expected.
\end{quote}

This defintion is derived, in part, from \cite{smithPolynyasLeadsOverview1990} and requires a
further, modern qualification, and that is that polynyas are \textit{non-linear}
regions \citep{massomDistributionFormativeProcesses1998} of open water and/or thin ice. The inverse of this
given by \cite{smithPolynyasLeadsOverview1990} in their definition of oceanographic \textit{leads} as
\textit{linear} openings in sea ice differentiated by their absence of
re-occurrence at a particular geographic location. For completeness, it should
be added, that coastal ocean polynyas are distinctly different from open ocean
polynyas in that the former are associated with sea ice production (\sip),
salt-rejection and hence \dsw formation
\citep{massomDistributionFormativeProcesses1998,maquedam.PolynyaDynamicsReview,tamuraMappingSeaIce2008}. This is in contrast to open ocean
polynyas (i.e. sensible heat polynyas) that are not typically associated with
\sip, especially in the Southern Ocean. Categorically, it is the water mass
modification occurring in coastal polynyas of Antarctica which are a
key driver of the \moc \citep{gawarkiewiczNumericalStudyDense1995,orsiCirculationMixingProduction1999,galton-fenziModelingBasalMelting2012,sansivieroModellingSeaIce2017},
and particularly, of concern in this work, is that the \cdp has been
shown to contribute 13\textendash 30\% of the total \aabw mass \citep{ohshimaAntarcticBottomWater2013}.

Given the identification of the \cdp as a significant contributor to \aabw, the
effort to determine its controlling factors has been pursued \citep{nakayamaNumericalInvestigationFormation2014,liuModelingModifiedCircumpolar2017}. \cite{fraserLandfastIceControls2019}
have shown that the contribution, specifically the length of landfast
ice (\fastice), termed Cape Darnley \fastice (\cdfi) is the most corollary physical feature (measured by satellite) in relation to the size of \cdp. (Note, I
include here the `dagger' portion, defined by \cite{fraserLandfastIceControls2019}, when
considering \cdfi). Beyond the physical importance of \fastice with polynya
dynamics it should be further ascribed that their role within the environment as unique habitats
is critical and our knowledge is formative \citep{trathanFirstRecordedLoss2011,ainleyApparentPopulationDecrease2015a}. It may also be worth mentioning here, that conceptually
differentiating \fastice from other types of sea ice (e.g. pack
ice, floes, etc.), I do so inline with common practice via the World
Meteorological Definition (\wmo) definition \cite{wmoWMOSeaiceNomenclature1970}:
\begin{quote}
  Sea ice which forms and remains fast along the coast, where it is attached to
  the shore, to an ice wall, to an ice front, between shoals or grounded
  icebergs. Vertical fluctuations may be observed during changes of sea-level.
  Fast ice may be formed in situ from sea water or by freezing of floating ice
  of any age to the shore, and it may extend a few metres or several hundred
  kilometres from the coast. Fast ice may be more than one year old and may then
  be prefixed with the appropriate age category (old, second-year, or multi-
  year).
\end{quote}

In many parts of the Antarctic continental shelf bathymetry favours the grounding of
icebergs \citep{wescheIcebergSignaturesDetection2012}. The grounding of icebergs provides a linking/trapping/blocking
mechanism for sea ice to become fastened to the iceberg \citep{massomEffectsRegionalFastice2001,fraserLandfastIceControls2019,liSpatiotemporalPatternsLandfast2020}. In some cases this
fastening links back to the shoreline, shelf edge, grounding keels, glacial
tongues, or ice walls \citep{massomDistributionFormativeProcesses1998,fraserEastAntarcticLandfast2012}.
\cdfi (and hence \cdp) appear connected in this manner to the seafloor topography
adjacent to Cape Darnley (i.e. Fram Bank) \citep{nakayamaNumericalInvestigationFormation2014,liuModelingModifiedCircumpolar2017}. It has
been shown that the spatial and temporal distribution of \fastice acts as
barrier and either inhibits or an enhances the size a coastal ocean
polynya
\citep{massomEffectsRegionalFastice2001,ohshimaAntarcticBottomWater2013,nakayamaNumericalInvestigationFormation2014,nihashiCircumpolarMappingAntarctic2015,kusaharaDenseShelfWater2017,fraserLandfastIceControls2019}.
The work done recently by \cite{fraserLandfastIceControls2019} provides a new metric for
\fastice categorisation and has shown that quantifying and predicting
\sip cannot rely solely on representation of the near surface wind, but also
must include the spatial and temporal distribution of the \fastice. Hence, a
growing \fastice database of Antarctic is underway \citep{fraserHighresolutionMappingCircumAntarctic2020} and will
be utilised in this work.

Prior to the above work, a lack of appropriate repressentation of
\fastice as well as the critical iceberg distribution, at a given time, in a coupled
ice-ocean numerical model, under represents \sip, and hence \dsw formation,
incorrectly, or not at all \citep{nakayamaNumericalInvestigationFormation2014,liuModelingModifiedCircumpolar2017}. Fast\textendash  
ice estimates derived from satellite indicate the \cdp to be a region
of high \sip and \dsw formation (13\textendash 30\% of total \aabw mass) through annual
production of sea ice between 195 \(\pm\) 71 \si{\kilo\meter\cubed}
\citep{ohshimaAntarcticBottomWater2013,tamuraSeaIceProduction2016}. The
controlling mechanisms for the \cdp are yet to be well documented but
are likely both atmospheric and oceanographic via: katabatic surface
winds, westward circumpolar flow, and favourable bathymetric topology
of the subsurface northward extension of the \cdp. 

The impetus of this thesis is focused on correctly representing
\fastice via a coupled sea ice\textendash ocean model \textemdash the
Modular Ocean Model (MOM6), Los Alamos Community Sea Ice model (CICE)
\citep{hunkeSeaiceModelsClimate2010,hunkeCICEAlamosSea2010} and
Wavewatch III (WW3). Rationale of choosing these models as well as an
overview of oceanographic, atmospheric and cryospheric
features and concepts will now be provided.

\section{Oceanography of the Southern Ocean}
\label{sec:org3a4ad91}

The Southern Ocean is broadly defined as the body of planetary water
that separates Antarctica from the continental landmasses of South
America, Africa and Australia. The Antarctic Circumpolar Current
(\acc) is the major oceanographic feature that encompases the Southern
Ocean. An unimpeded flow of a series of ocean surface jets that
extend hundreds of metres into the water column, encircling the
planet and allowing for information in the form of heat, chemicals and
nutrients to be exchanged between the ocean basins on timescales of
months to years. It is the largest current of
water on earth, driven by westerly winds between 45-55 deg.S. and
buoyancy forcing that tilt surfaces of density (isopycnal surfaces) in
the ocean basins (Atlantic, Indian and Pacific) towards the south and create the eastward geostrophic
flow of the \acc. Waters of the ocean basins are thus connected via
the \acc\citep{orsiMeridionalExtentFronts1995}. This relationship gives rise to highly distinguishable meridional
gradient that separates the Southern Ocean waters from those of the
low-latitude ocean basin waters. The area where these waters converge
on the surface is the \acc or also referred to as the sub-tropical
front \citep{orsiCoolingVentilatingAbyssal2001}. In this manner it can be
thought of as a series of ocean jets that, due to the absence of any
pressure gradient force that a blocking land mass would impose, in
which no net geostrophic flow across the \acc can occur.

The \acc differs fundamentally from other major current
systems (western boundary currents) in that its zonal ocean transport
is known to be proportional \textit{directly} to the wind stress
\citep{munkNoteDynamicsAntarctic1951}. Whereas, western boundary
currents (Gulf Stream, Kuroshio, EAC, etc.), the wind stress is
balanced with lateral friction, and hence the transport is proportional to the
\textit{curl} of the wind stress
\citep{munkNoteDynamicsAntarctic1951}. \cite{hughesWhyWesternBoundary2001}
further showed that the interaction of the \acc flow with topographic features is also
fundamentally a balance in the absolute vorticity at depth and
advection of planetary vorticity, whereas the flat-bottom Sverdrup
balance between the curl of the wind stress and advection of planetary
vorticity, in the other ocean basins, is \textit{not} linked to bottom
topography. This fundamental distinction is
significant in that the transfer of energy from the atmosphere to
ocean, in the Southern Ocean, is not only like no where else on the
planet, but also intimately tied to the synoptic atmosphere of the
Southern Hemisphere and the bathymetry of the Southern Ocean.

Consider again the isopycnal surfaces of the ocean basins. In areas
where they rapidly rise in the southward direction, and are in balance
with geostrophic flow, they induce eastward flow intensification in
\acc. In areas where thesse sloping surfaces become baroclinically unstable meso-scale eddies are generated and facilitate
the transfer and the mixing of sub-polar and sub-tropical ocean basin waters. These
anti-cyclonic (warm core eddies) transport warm salty sub-tropical
nutrient poor waters towards the Antarctic shelf
\citep{rintoulGlobalInfluenceLocalized2018}. These warm core eddies,
with their divergent surface flow, facilitate the vertical momemtum
transport necessary to balance the heat loss of the sub-polar ocean
with that gained from the sub-tropical waters -- i.e. venting of the
deep ocean is ocean is achieved. Critical to this process is that the
vertical momemtum transport induces a deep water geostropic flow, in
turn balanced by the pressure gradient force over bathymetry, which
results in bottom force stress that opposes the surface wind. This
critically allows for transport where otherwise no time-averaged
geogstrophic flow can occur. Hence the \moc and the \acc are
initimately linked.

Along the coastline of Antarctica fierce katabatic oceanward winds
pour off the massive ice shelves and glaciers and have a slightly more
eastward tendency and hence drive Ekman transport of relatively cold
freshwater westward and away from the coast. This is observed as two
current systems called the Antarctic Slope Front and Antarctic Coastal
Current \citep{heywoodFateAntarcticSlope2004}. Distinguishing between
these two current systems is not necessarily beneficial and hence are
commonly referred to as the Antarctic Slope Current (\asc). Just as the \acc is a
series of eastward boundary jets in the deep water of the Southern Ocean, the
\asc is a series of westward jets that roughly follow the
continental shelf and control the rate at which waters flow onto and
away from the shelf waters \citep{thompsonAntarcticSlopeCurrent2018}. It is also important to
note that the \asc is not purely driven by local effects but also
linked to climatic systems El Nino Southern Oscillation (ENSO) and the
Southern Annular Mode (SAM) \citep{spenceRapidSubsurfaceWarming2014,nakayamaOriginCircumpolarDeep2018}. The \asc is
near surface current that flows at rates of 10-30 cm/s. The importance
of this 

\section{Cryosphere of Antarctic Coastal Margin}
\label{sec:org06d6cc3}

The continental shelf waters of Antarctica are the connector to the
Antarctic cryosphere. In West Antarctica this has been most pointedly
observed as the increase to the heat
content of the shelf waters has directly contributed to dramatic
thinning of those ice shelves [Heywood2014,Paolo2015]. The thinning
has been directly attributed to melting from beneath the ice tongues
and shelves as the warmer water that has arrived from the sub-tropical
basins via a set of dynamic processes and provides the annual heat required
to thin these massive quantities of ice [Pritchard 2012]. The present rate at
which thinning/melting is occurring is three times what it was three
decades ago [Rintoooul2016].

The depth of the continental shelf around most of Antarctica is 500m
with many important high gradient canyons crucial in hydrographic
feature definition of the cryosphere. 

\section{Coupled Sea Ice -- Ocean Modelling}
\label{sec:orgd1f8a57}
\subsection{MOM6}
\label{sec:orgdcd8ebd}
\subsection{CICE6}
\label{sec:org37c7cd8}
\subsection{WW3}
\label{sec:orgb753653}
\section{References}
\label{sec:orge3e57fe}
\printbibliography[heading=none]
\end{document}